Este anexo contiene un listado de funcionalidades. Se dividen en secciones en función de la entidad a la que "pertenecen", con la que más estén relacionadas, o la que afecten principalmente. Cada funcionalidad tiene:
\begin{itemize}
    \item Identificador con el formato REQ-X, donde X es un número único a cada funcionalidad
    \item Nivel: 0 para funcionalidades básicas, 1 para funcionalidades avanzadas esenciales, 2 para funcionalidades avanzadas opcionales
    \item Descripción de la funcionalidad
\end{itemize}
Respecto a qué tipo de usuario puede acceder a qué funcionalidades, se describe a continuación:
\begin{itemize}
    \item \textbf{Usuario Anónimo}: únicamente puede acceder a las funcionalidades REQ-0 y REQ-3, creación de cuenta e inicio de sesión
    \item \textbf{Usuario Autenticado}: puede acceder a todas las demás funcionalidades. Las funcionalidades de edición y eliminación de una entidad tienen como prerrequisito que el usuario sea propietario.
    \item \textbf{Usuario Administrador}: puede acceder las mismas funcionalidades que el Usuario Autenticado, y puede eliminar recursos de los que no es propietario (actividades, rutas, murales) únicamente cuando estos NO tengan visibilidad privada. Esto no se extiende a la modificación; el usuario administrador únicamente puede modificar recursos de los que sea propietario.
\end{itemize}

\section{Usuario}

\begin{longtable}{p{2cm} p{1cm} p{11cm}}
\toprule
\textbf{REQ-X} & \textbf{Nivel} & \textbf{Descripción} \\
\midrule
\endhead

REQ-1 & 0 & Crear cuenta \\
REQ-2 & 0 & Eliminar cuenta \\
REQ-3 & 0 & Iniciar sesión \\
REQ-4 & 0 & Cerrar sesión \\
REQ-5 & 0 & Visualizar información de una cuenta (Nombre de usuario, Nombre público, Correo electrónico) \\
REQ-6 & 0 & Editar datos de una cuenta \\
REQ-7 & 0 & Visualizar datos sobre las actividades y rutas creadas:
\begin{itemize}
    \item Panel de información general (Nombre de usuario, Número de rutas y actividades, Distancia total, Tiempo total de entrenamiento, Ritmo medio,  Máxima distancia en una actividad, Máximo tiempo en una actividad)
    \item Panel de actividades recientes
    \item Panel de rutas creadas
\end{itemize} \\
REQ-8 & 1 & Visualizar récords:
\begin{itemize}
    \item Mínimo tiempo en recorrer distancias fijas (1km, 5km, 10km, 21km, 42km)
    \item Máxima distancia recorrida en periodos fijos (1, 5, 10, 30, 60 minutos)
\end{itemize} \\

\bottomrule
\end{longtable}

\section{Actividad}
\subsection{Individual}

\begin{longtable}{p{2cm} p{1cm} p{11cm}}
\toprule
\textbf{REQ-X} & \textbf{Nivel} & \textbf{Descripción} \\
\midrule
\endhead

REQ-9 & 0 & Importar una actividad en formato GPX \\
REQ-10 & 0 & Eliminar una actividad \\
REQ-11 & 0 & Editar visibilidad de una actividad \\
REQ-12 & 0 & Editar el nombre de una actividad \\
REQ-13 & 0 & Cambiar la ruta asignada a una actividad \\
REQ-14 & 0 & Visualizar un resumen de la actividad (Nombre, Fecha y hora de inicio, Duración, Distancia, Elevación ganada/perdida, Ritmo medio) \\
REQ-15 & 0 & Visualizar métricas en un gráfico de líneas (Altitud, Ritmo) \\
REQ-16 & 2 & Cambiar el valor del eje X (Tiempo transcurrido, Distancia) \\
REQ-17 & 0 & Visualizar métricas en histogramas (Altitud, Ritmo) \\
REQ-18 & 2 & Modificar el número de secciones del histograma \\
REQ-19 & 0 & Visualizar percentiles, media, IQR, porcentaje de valores anómalos y desviación media \\
REQ-20 & 0 & Establecer valor objetivo, límite superior e inferior para una métrica \\
REQ-21 & 0 & Visualizar el porcentaje de valores en cada área definida \\
REQ-22 & 0 & Visualizar límites y valor objetivo en el gráfico \\
REQ-23 & 0 & Visualizar la actividad en un mapa \\
REQ-24 & 0 & Estudiar métricas:
(Pulsaciones, Cadencia, Ritmo) \\
REQ-25 & 2 & Miniatura del mapa al hacer \textit{hover} sobre listas \\

\bottomrule
\end{longtable}

\subsection{Comparación}

\begin{longtable}{p{2cm} p{1cm} p{11cm}}
\toprule
\textbf{REQ-X} & \textbf{Nivel} & \textbf{Descripción} \\
\midrule
\endhead

REQ-26 & 1 & Visualizar resumen comparativo de métricas (distancia, tiempo, elevación, ritmo) \\
REQ-27 & 1 & Visualizar tabla comparativa con métricas estadísticas \\
REQ-28 & 1 & Calcular porcentaje de valores entre dos límites definidos por el usuario \\
REQ-29 & 1 & Visualizar tabla de récords de tiempo y distancia \\
REQ-30 & 2 & Visualizar divisiones de distancia configurables \\
REQ-31 & 1 & Visualizar métricas en gráficos de líneas para varias actividades \\
REQ-32 & 2 & Visualizar métricas en histogramas para varias actividades \\
REQ-33 & 1 & Visualizar varias actividades en un mismo mapa \\
REQ-34 & 1 & Visualizar varias actividades en mapas distintos \\
REQ-35 & 2 & Cambiar qué actividades se muestran \\
REQ-36 & 2 & Cambiar estrategia de resolución de duplicados \\

\bottomrule
\end{longtable}

\section{Ruta}

\begin{longtable}{p{2cm} p{1cm} p{11cm}}
\toprule
\textbf{REQ-X} & \textbf{Nivel} & \textbf{Descripción} \\
\midrule
\endhead

REQ-37 & 0 & Crear una ruta a partir de una actividad \\
REQ-38 & 0 & Eliminar una ruta \\
REQ-39 & 0 & Añadir actividad a una ruta \\
REQ-40 & 0 & Eliminar actividad de una ruta \\
REQ-41 & 1 & Comparar actividades en una ruta \\
REQ-42 & 0 & Editar visibilidad de una ruta \\
REQ-43 & 2 & Filtrar rutas por visibilidad \\
REQ-44 & 0 & Visualizar listado de rutas \\
REQ-45 & 0 & Visualizar actividades de una ruta \\
REQ-46 & 0 & Visualizar información de una ruta:
(Nombre, Distancia total, Descripción, Elevación ganada/perdida, Mapa) \\
REQ-47 & 2 & Visualizar datos acumulativos \\
REQ-48 & 2 & Visualizar mejor actividad \\
REQ-49 & 2 & Visualizar récords \\

\bottomrule
\end{longtable}

\section{Mural}

\begin{longtable}{p{2cm} p{1cm} p{11cm}}
\toprule
\textbf{REQ-X} & \textbf{Nivel} & \textbf{Descripción} \\
\midrule
\endhead

REQ-50 & 0 & Crear mural \\
REQ-51 & 0 & Borrar mural \\
REQ-52 & 0 & Editar mural \\
REQ-53 & 2 & Traspasar propiedad \\
REQ-54 & 0 & Unirse a un mural mediante código \\
REQ-55 & 0 & Expulsar usuario \\
REQ-56 & 1 & Banear usuario \\
REQ-57 & 0 & Visualizar y comparar actividades visibles para el mural \\
REQ-58 & 0 & Visualizar rutas visibles para el mural \\
REQ-59 & 0 & Visualizar dashboard del mural con panel de miembros \\
REQ-60 & 0 & Listar murales creados \\
REQ-61 & 0 & Listar murales en los que se es miembro \\
REQ-62 & 0 & Listar murales disponibles \\
REQ-63 & 2 & Invitar usuarios a un mural \\
REQ-64 & 0 & Abandonar mural \\

\bottomrule
\end{longtable}